\chapter{Calcul différentiel et sous-variétés différentiables de $\Rn$}

\section{Sous-variétés de $\Rn$}

\begin{definition}[Sous-variété] \label{def: Sous-variété}
	Un sous-ensemble $M \subset \Rn$ est une \textbf{sous-variété} de dimension $m \leqslant n$ et de classe $C^k$ si pour tout $p \in M$, il existe un ouvert $U \in \Rn$ qui contient $p$ et un difféomorphisme $\phi: U \to V$ de classe $C^k$ tel que
	\[
	\phi(U \cap M) = V \cap E
	\]
	où $E \subset \Rn$ est un sous-espace affine de dimension $m$. On dira que $(n-m)$ est la \textbf{codimension} de $M$.
\end{definition}

\begin{remarque}
	Sans perte de généralité, quitte à composer le difféomorphisme $\phi$ avec une application affine, on peut supposer $E$ est une sous-espace vectoriel. De plus on peut aussi supposer que $E$ est de la forme
	\[
	E = \{ x \in \Rn \mid x_{m+1} = \ldots = x_n = 0\}
	\]
	Et donc la pratique, vérifier qu'un sous-ensemble $M$ est un sous-variété de dimension $m$ revient à montrer que pour tout $p \in M$, il existe un \emph{système coordonnées curvilignes} $y_1, \ldots, y_n$ dans un voisinage $U \subset \Rn$ de $p$ tel que
	\[
	\phi(U \cap M) = \{y \in V \mid y_{m+1} = \ldots = y_n = 0\}
	\]
	où $\phi$ désigne le difféomorphisme qui représente le changement de coordonnées.
\end{remarque}

\begin{theoreme}
	Soit $f : U  \subset \R^m \to \Rn$ une application de classe $C^k$ et de rang constant $r$. Alors on a
	\begin{enumerate}
		\item Pour tout $q \in \Rn$, la préimage
		\[
		f^{-1} (q) \subset U
		\]
		est une sous-variété de codimension $r$ dans $\Rm$
		
		\item Pour tout $p \in U$, il existe un voisinage $V_p \subset U$ tel que l'image
		\[
		f(V_p) \subset \Rn
		\]
		est une sous-variété de dimension $r$.
	\end{enumerate}
\end{theoreme}